% Plantilla base de LEMA. Cada tema la carrega amb \input després de \documentclass.
% Només fa servir paquets que ja hi ha a la instal·lació de MiKTeX de l'usuari.
\usepackage[T1]{fontenc}
\usepackage[utf8]{inputenc}
\usepackage[catalan]{babel}
\usepackage{lmodern}
\usepackage{amsmath,amssymb,amsthm}
\usepackage{xcolor}

% Pàgina de llibre amb un marge dret ample per a les fonts i les notetes.
\setlength{\textwidth}{12.4cm}
\setlength{\textheight}{22.5cm}
\setlength{\oddsidemargin}{0.3cm}
\setlength{\evensidemargin}{0.3cm}
\setlength{\topmargin}{-0.6cm}
\setlength{\marginparsep}{0.5cm}
\setlength{\marginparwidth}{3.7cm}
\setlength{\parindent}{0pt}
\setlength{\parskip}{0.4\baselineskip}

% Entorns de cada tipus de bloc. Numeració comuna dins del tema (1.1, 1.2...).
\theoremstyle{plain}
\newtheorem{teorema}{Teorema}[chapter]
\newtheorem{proposicio}[teorema]{Proposició}
\newtheorem{lema}[teorema]{Lema}
\newtheorem{corollari}[teorema]{Corol·lari}
\theoremstyle{definition}
\newtheorem{definicio}[teorema]{Definició}
\newtheorem{exemple}[teorema]{Exemple}
\newtheorem{observacio}[teorema]{Observació}
\newtheorem{notacio}[teorema]{Notació}
\newtheorem{exercici}[teorema]{Exercici}
% Llistes numerades com als apunts: 1), 2)...
\renewcommand{\labelenumi}{\theenumi)}
\addto\captionscatalan{\renewcommand{\proofname}{Demostració}}

% Colors de les marques.
\definecolor{lemacorr}{HTML}{B4231B}   % canvi respecte als apunts
\definecolor{lemanota}{HTML}{1F4FB8}   % notes i notetes
\definecolor{lemapend}{HTML}{B45309}   % dubtes i coses pendents
\definecolor{lemafont}{HTML}{6B7280}   % d'on ve cada bloc
\definecolor{lemafons}{HTML}{F1F4F9}

% D'on ve el bloc: font i pàgina dels apunts, al marge.
\newcommand{\origen}[1]{\marginpar{\raggedright\footnotesize\color{lemafont}#1}}
% Il·legible o dubtós.
\newcommand{\dubte}[1]{\textcolor{lemapend}{[#1?]}}
% Canvi respecte als apunts.
\newcommand{\correccio}[1]{\textcolor{lemacorr}{#1}\marginpar{\raggedright\footnotesize\color{lemacorr}Correcció respecte als apunts}}
% Discrepància vàlida amb la font.
\newcommand{\nota}[1]{\marginpar{\raggedright\footnotesize\color{lemanota}#1}}
% Noteta numerada: el número al text i al marge; la llista sencera és a l'informe.
\newcommand{\noteta}[2]{\textsuperscript{\textcolor{lemanota}{\sffamily\bfseries #1}}\marginpar{\raggedright\footnotesize\color{lemanota}\textsf{\textbf{#1}} #2}}
% Caixa per al que no es va veure a classe i s'ha pres d'una font.
\newcommand{\noclasse}[2]{\par\smallskip\noindent\fcolorbox{lemanota}{lemafons}{\parbox{\dimexpr\linewidth-2\fboxsep-2\fboxrule\relax}{\footnotesize\textbf{No vist a classe}, tret de #1.\par\normalsize #2}}\par\smallskip}
% Demostració que encara no tenim.
\newcommand{\pendent}[1]{\par\noindent\textit{Demostració.}\ \textcolor{lemapend}{\textbf{Pendent:} #1}\par}
